\documentclass[12pt, a4paper, oneside]{book}
\usepackage{fontspec}
%\setmainfont{Arial}[ItalicFont={Arial Italic}]
%\setmainfont{Gill Sans MT}[ItalicFont={Gill Sans MT Italic}]
\usepackage[utf8]{inputenc}
\usepackage[margin=1.5cm, bindingoffset=1cm]{geometry}
\linespread{1.5}
\usepackage{float}
\usepackage{multicol}
\usepackage{csquotes}
\usepackage{subfig}
\usepackage{graphicx}
\usepackage{wrapfig}
\usepackage[table]{xcolor}
\usepackage{indentfirst}
\setlength{\parindent}{0cm}
\usepackage[italian]{babel}
\usepackage{amsmath,amssymb}
\usepackage{hyperref}
\usepackage{color}
\usepackage{listings}
\usepackage{wrapfig}
\usepackage{url}
\lstset{showstringspaces=false}

%TABELLE
\usepackage{booktabs} % Per le linee orizzontali di alta qualità
\usepackage{tabularx} % Per tabelle di larghezza fissa con colonne espandibili
\definecolor{headercolor}{gray}{0.9} % Imposta il colore dell'intestazione a grigio
\usepackage{array}
\setlength{\arrayrulewidth}{0.3mm}

\usepackage{eurosym} %simbolo Euro

\usepackage{fancyhdr}
\pagestyle{fancy}
\renewcommand{\chaptermark}[1]{%
    \markboth{\chaptername
    \ \thechapter.\ #1}{}}
\fancyhf{}
\fancyhead[L]{\textsl{\leftmark}}
\fancyhead[R]{\textsl{IS- Problem Statement}}
\fancyfoot[C]{\thepage}


\title{Problem Statement - IS}
\author{Chiara Coscarelli, Angelo Fiorillo, Marco Santoriello, Fabio Sessa}
\date{Ottobre 2023}

\begin{document}
    \begin{titlepage}
        \begin{center}
            \LARGE{\uppercase{Università degli Studi di Salerno}}\\
            \vspace{5mm}
            %Dipartimento
            \uppercase{\normalsize Dipartimento di Informatica}\\
        \end{center}
        \begin{figure}[H]
            \centering
            \includegraphics[width=0.3\textwidth]{img/logo_unisa.png}
        \end{figure}

        \begin{center}
            %Corso di Laurea
            \normalsize{ Corso di Laurea in informatica }\\
            \vspace{5mm}
            %Titolo
            {\LARGE{\bf Problem Statement - iSnapGaming }}\\
            \vspace{10mm}
        \end{center}

        \vspace{15mm}
        \noindent
        \begin{minipage}[t]{0.47\textwidth}
        %Relatore
        {\large{ Docente:\\\bf  Prof. Andrea De Lucia}}
            \vspace{12mm}\\
        \end{minipage}
        \hfill
        \begin{minipage}[t]{0.4\textwidth}\raggedright
        %Candidato
        {\large{Componenti: \\ \bf Chiara Coscarelli\\Angelo Fiorillo\\Marco Santoriello\\Fabio Sessa}}
        \end{minipage}

        \vspace{20mm}

        \vspace{50mm}
        %Anno Accademico
        \centering{\large \uppercase{ Anno Accademico 2023/2024 }}

    \end{titlepage}
    \newpage
%MEMBRI DEL PROGETTO
    \begin{center}
    {\LARGE{Membri del progetto}}
    \end{center}
    \begin{center}
        \begin{tabular}{|c|c|m{8cm}|}
            \hline
            \cellcolor{lightgray} \textbf{Nome} & \cellcolor{lightgray} \textbf{Matricola} & \cellcolor{lightgray} \textbf{E-mail}\\ \hline
            Coscarelli  Chiara & 0512113869 & c.coscarelli1@studenti.unisa.it\\ \hline
            Fiorillo Angelo & 0512105739 & a.fiorillo9@studenti.unisa.it\\ \hline
            Santoriello Marco &  0512114100 & m.santoriello37@studenti.unisa.it\\ \hline
            Sessa Fabio & 0512114634 & f.sessa30@studenti.unisa.it\\ \hline
        \end{tabular}
    \end{center}

    \begin{center}
    {\LARGE{Revision History}}
    \end{center}

    %TABELLA REVISION HISTORY
    \begin{center}
        \begin{tabular}{|c|c|m{8cm}|m{3.5cm}|}
            \hline
            \cellcolor{lightgray} \textbf{Data} & \cellcolor{lightgray} \textbf{Versione} & \cellcolor{lightgray} \textbf{Descrizione} & \cellcolor{lightgray} \textbf{Autore}\\ \hline
            09/10/2023 & 1.00 & Impostazione documento e descrizione del problema & Marco Santoriello\\ \hline
            10/10/2023 & 1.01 & Aggiunta degli scenari Sxx, Sxx & Marco Santoriello\\ \hline
            11/10/2023 & 1.02 & Aggiunta di alcuni scenari & Fabio Sessa\\ \hline
            11/10/2023 & 1.03 & Aggiunta di alcuni scenari & Chiara Coscarelli\\ \hline
            11/10/2023 & 1.04 & Aggiunta scenario ricerca prodotti & Angelo Fiorillo\\ \hline
            15/10/2023 & 1.05 & Aggiunta scenario gestione account e stato ordine & Angelo Fiorillo\\ \hline
        \end{tabular}
    \end{center}





    \setcounter{tocdepth}{3} %IMPOSTO LIVELLO PROFONDITA' INDICE

    \tableofcontents

    \chapter*{Problem Statement}
%La seguente istruzione aggiunge il capitolo all'indice, siccome il comando
%con la star * non inserisce il numero del capitolo/sezione ma non mostra
%%nemmeno la sezione nell'indice
    \addcontentsline{toc}{chapter}{\protect\numberline{}Problem Statement}
    \section*{Problem}
    \addcontentsline{toc}{section}{\protect\numberline{}Problem}
    iSnapGaming nasce come un e-commerce specializzato nella vendita di videogiochi. Non ha una sede fisica, dunque l'unica entrata del business è rappresentata esclusivamente dalla piattaforma di vendite online.

    \section*{Scenarios}
    \addcontentsline{toc}{section}{\protect\numberline{}Scenarios}
    \subsection*{Sxx: Acquisto cliente loggato}
    \addcontentsline{toc}{subsection}{\protect\numberline{}Sxx: Acquisto cliente loggato}
    L'utente \textit{Mario Rossi} accede alla pagina principale della piattaforma, clicca sul pulsante per accedere al suo account e viene reindirizzato alla pagina di \textit{login}. A questo punto, compila i campi del form \textit{username} e \textit{password} con i propri dati (username: marrossi, password: MyPassWord11) e lo conferma cliccando sull'opzione preposta. A questo punto, viene rimandato sulla home page, dalla quale sceglie un prodotto tra quelli in evidenza e lo aggiunge al carrello.\\
    Clicca successivamente sulla funzionalità del carrello e viene portato sulla pagina del carrello, dove \textit{incrementa} la quantità del prodotto da acquistare a 2.\\
    Sicuro del prodotto, della quantità che intende acquistare ed del prezzo finale, pari a \officialeuro 55, seleziona l'opzione di \textit{check-out}, che lo rimanda alla pagina di inserimento dei dati di spedizione e del metodo di pagamento. Compila i vari campi: \textit{Via: Via Carmine}, \textit{Civico: 57}, \textit{Città: Salerno}, \textit{Provincia: SA}, \textit{Stato: Italia}, \textit{Cellulare: 345 670 9276}, \textit{Numero di carta: 5333547887933563}, \textit{Scadenza: 10/25}, \textit{CVV: 679}. Infine, seleziona l'opzione per terminare l'acquisto.\\
    L'ordine viene acquisito dal sistema, che, verificati i dati relativi al pagamento, svuota il carrello dell'utente, il quale viene reindirizzato alla \textit{thank-you-page}, dove gli verrà mostrato il riepilogo dell'ordine, con il relativo codice.

    \subsection*{Sxx: Acquisto cliente non loggato}
    \addcontentsline{toc}{subsection}{\protect\numberline{}Sxx: Acquisto cliente non loggato}
    L'utente \textit{Mario Rossi} accede alla pagina principale della piattaforma e si porta, attraverso l'apposita opzione presente nella barra di navigazione, alla pagina del catalogo dalla quale sceglie due videogiochi e li aggiunge al carrello. Si reca, dunque, alla pagina del carrello, dove gli vengono mostrati i due prodotti presenti nel carrello, le opzioni per \textit{modificare la quantità} di ciascun prodotto e il \textit{totale}, che ammonta a \officialeuro 90,00. A questo punto, seleziona la funzionalità per effettuare il check-out e, non essendo loggato, viene reindirizzato alla pagina di \textit{login}. Una volta effettuato l'accesso, viene riportato sulla \textit{Home page}. Da qui, accede nuovamente al carrello e seleziona l'opzione di \textit{check-out}, che lo rimanda alla pagina di inserimento dei dati di spedizione e del metodo di pagamento. Compila i vari campi: \textit{Via: Via Carmine}, \textit{Civico: 57}, \textit{Città: Salerno}, \textit{Provincia: SA}, \textit{Stato: Italia}, \textit{Cellulare: 345 670 9276}, \textit{Numero di carta: 5333547887933563}, \textit{Scadenza: 10/25}, \textit{CVV: 679}. Infine, seleziona l'opzione per terminare l'acquisto.\\
    L'ordine viene acquisito dal sistema, che, verificati i dati relativi al pagamento, svuota il carrello dell'utente, il quale viene reindirizzato alla \textit{thank-you-page}, dove gli verrà mostrato il riepilogo dell'ordine, con il relativo codice.

    \subsection*{Sxx: Acquisto cliente non registrato}
    \addcontentsline{toc}{subsection}{\protect\numberline{}Sxx: Acquisto cliente non registrato}
    L'utente Luca Neri accede alla piattaforma dal suo browser per la prima volta, avendone sentito parlare da alcuni conoscenti, entrando nella pagina principale. In particolare, rimane colpito da un videogioco presente tra i \textit{featured-products} e decide di acquistarlo. Dunque, per prima cosa, aggiunge il prodotto al carrello e si reca sulla pagina dedicata al carrello. Seleziona l'opzione di \textit{check-out}, ma, non avendo ancora un account, il sito lo reindirizza alla pagina di registrazione, dove compila i vari campi necessari alla registrazione: \textit{Nome: Luca}, \textit{Cognome: Neri}, \textit{Data di nascita: 27/05/1998}, \textit{Codice Fiscale: LCUNRE98E27F205X}, \textit{Username: nlucas98}, \textit{Indirizzo e-mail: lneri98@gmail.com} e \textit{Password: LaMiaPasswordSafe12}. Infine, clicca su \textit{Registrati}. La registrazione termina con successo, e l'utente viene riportato sulla \textit{Home page}. Da qui, accede nuovamente alla pagina del carrello, prosegue al check-out e viene rimandato alla schermata per ultimare il pagamento.\\
    Compila i vari campi: \textit{Via: Via Piave}, \textit{Civico: 21}, \textit{Città: Voghera}, \textit{Provincia: PV}, \textit{Stato: Italia}, \textit{Cellulare: 327 352 3770}, \textit{Numero di carta: 5333756908563312}, \textit{Scadenza: 12/24}, \textit{CVV: 550}. Infine, seleziona l'opzione per terminare l'acquisto.\\
    L'ordine viene acquisito dal sistema, che, verificati i dati relativi al pagamento, svuota il carrello dell'utente, il quale viene reindirizzato alla \textit{thank-you-page}, dove gli verrà mostrato il riepilogo dell'ordine, con il relativo codice.

    \subsection*{Sxx: Modifica di un prodotto}
    \addcontentsline{toc}{subsection}{\protect\numberline{}Sxx: Modifica di un prodotto}
    Il \textit{gestore dei prodotti} Michele Verdi entra nella piattaforma, si autentica attraverso la schermata di login ed accede alla pagina del catalogo. Da qui, entra nella pagina del videogioco \textit{Assassin's Creed: Valhalla}, e clicca sull'opzione che consente di modificare il prodotto, dato che vuole farlo apparire come \textit{non disponibile}, a causa del ritardo del fornitore che avrebbe dovuto consegnargli le nuove copie. Seleziona, quindi, il campo \textit{Disponibile}, e lo imposta su \textit{No}. Clicca sul pulsante di salvataggio delle modifiche, e viene reindirizzato alla pagina del prodotto, che presenta la modifica appena effettuata.

    \subsection*{Sxx: Aggiunta di un nuovo gestore degli ordini}
    \addcontentsline{toc}{subsection}{\protect\numberline{}Sxx: Aggiunta di un nuovo gestore degli ordini}
    L'\textit{admin}, cioè l'amministratore dell'intero sistema, Francesco Bianchi, accede al suo account e viene rimandato sulla HomePage. Da qui, seleziona l'opzione \textit{Aggungi Nuovo} da un menu dedicato. Gli compare la pagina che permette di inserire un nuovo utente. Compila i vari campi: \textit{Nome: Antonio}, \textit{Cognome: Romano}, \textit{Data di nascita: 15/03/2000}, \textit{Codice Fiscale: RMNNTN00C15H703B}, \textit{Username: goRomano}, \textit{Indirizzo e-mail: aRomano27@gmail.com} e \textit{Password: LaTuaPassword1} e, infine, seleziona la funzionalità che permette di confermare l'aggiunta.

    \subsection*{Sxx: Aggiunta di un prodotto}
    \addcontentsline{toc}{subsection}{\protect\numberline{}Sxx: Aggiunta di un prodotto}
    Il gestore dei prodotti Marco Bianchi deve aggiungere il gioco Cyberpunk 2077 al catalogo. Il gestore
    effettua l'accesso selezionando la sezione di autenticazione dove inserisce i dati: Email: mbianchi02@gmail.com, Password: Mari0bianchi
    Il sistema avverte il gestore che l'autenticazione è andata a buon fine. Il gestore accede alla sezione
    Aggiungi Prodotto che lo porta alla schermata dedicata. Il gestore inserisce le seguenti informazioni:
    ID: x123P, Nome: Cyberpunk 2077, Produttore: CD Projekt RED, Piattaforma: PS5, Prezzo: 49.99, Quantità: 250,
    Categoria: Action RPG, PEGI: 18, Anno di Produtzione: 2020. Il gestore invia le informazioni
    e il sistema avverte che l'operazione di aggiunta è stata eseguita con successo.

    \subsection*{Sxx: Rimozione di un prodotto}
    \addcontentsline{toc}{subsection}{\protect\numberline{}Sxx: Rimozione di un prodotto}
    Il gestore dei prodotti Marco Bianchi deve rimuovere il gioco Cyberpunk 2077 dal catalogo.
    Il gestore effettua l'accesso selezionando la sezione di autenticazione dove inserisce i dati: Email: mbianchi02@gmail.com, Password: Mari0bianchi
    Il sistema avverte il gestore che l'autenticazione è andata a buon fine. Il gestore accede alla sezione del Catalogo.
    Il gestore utilizza la funzionalità della ricerca per categoria dei prodotti e seleziona "Action RPG". Il gestore
    trova Cyberpunk 2077 e accede all'area riservata ai dettagli. Il gestore seleziona la funzionalità che lo
    permette di eliminare definitivamente il prodotto dal catalogo. Il sistema avverte il gestore che l'operazione è avvenuta con successo.

    \subsection*{Sxx: Aggiunta di una recensione}
    \addcontentsline{toc}{subsection}{\protect\numberline{}Sxx: Aggiunta di una recensione}
    L'utente Marco Marroni ha un obiettivo ben preciso: recensire il gioco Assassin's Creed Mirage che ha appena acquistato.
    Prima di tutto, si autentica al sistema inserendo le proprie credenziali, ovvero Email: mmarroni02@gmail.com, Password: Marc0marroni
    Il sistema avverte l'utente che l'autenticazione è andata a buon fine, dato che ha inserito i dati correttamente.
    L'utente accede all'area riservata del proprio account e successivamente utilizza la funzionalità
    che lo porta nella schermata dell'elenco degli ordini effettuati. La schermata presenta tutti gli ordini,
    da quelli completati a quelli ancora in spedizione. L'utente trova l'ordine in cui ha acquistato
    Assassin's Creed Mirage e accede alla sezione che lo permette di inserire la propria recensione
    in un campo di testo apposito, digitando il seguente messaggio: "Il prodotto è arrivato in condizioni perfette
    e lo consiglio a tutti gli appassionati!". L'utente utilizza la funzionalità che permette l'invio della recensione.
    Successivamente, il sistema avverte l'utente di attendere che la pubblicazione della sua recensione sia approvata
    da un gestore delle recensioni.

    \subsection*{Sxx: Approvazione di una recensione}
    \addcontentsline{toc}{subsection}{\protect\numberline{}Sxx: Pubblicazione di una recensione}
    Il gestore delle recensioni Luca Verdi, in una normale giornata di lavoro, deve controllare le varie richieste
    di recensioni (effettuate dagli utenti) che deve approvare e rendere pubbliche.
    Prima di tutto, il gestore si autentica al sistema inserendo le proprie credenziali, ovvero Email: lverdi02@gmail.com, Password: Lucaverd1
    Il sistema avverte il gestore che l'autenticazione va a buon fine. Dopodiché, il gestore accede alla schermata
    Gestione Recensioni per visualizzare l'elenco di tutte le richieste di recensione. Il gestore nota che in cima
    alla lista è presente una recensione dell'utente Marco Marroni riferita al gioco Assassin's Creed Mirage,
    con il seguente messaggio: "Il prodotto è arrivato in condizioni perfette e lo consiglio a tutti gli appassionati!".
    Il gestore capisce che il testo è privo di contenuti offensivi e decide di approvare la pubblicazione della recensione.
    Il sistema avverte il gestore che la pubblicazione è avvenuta con successo.

    *{Sxx: Registrazione}
    \addcontentsline{toc}{subsection}{\protect\numberline{}Sxx: Registrazione}
     L'utente Luca Neri, una volta entrato sulla piattaforma, accede al suo account utilizzando l'apposita funzionalità di "login", che lo
     reindirizza alla pagina di login del sito.
     Una volta qui, l'utente seleziona "Registrati Ora". Fatto ciò, viene immediatamente reindirizzato alla pagina di registrazione.
     Su questa pagina, Luca compila attentamente tutti i campi necessari. Nei campi obbligatori, inserisce le seguenti informazioni:
     Nome: Luca, Cognome: Neri, Data di nascita: 27/05/1998, Codice Fiscale: LCUNRE98E27F205X,
     Username: nlucas98, Indirizzo e-mail: lneri98@gmail.com e Password: LaMiaPasswordSafe12.
     La registrazione termina con successo, e l’utente viene riportato sulla Home page.

    \subsection*{Sxx: Login Cliente}
    \addcontentsline{toc}{subsection}{\protect\numberline{}Sxx: Login Cliente}
    Il cliente Mario Rossi, che si trova nella pagina del catalogo, accede alla pagina di Login selezionando l'apposita funzionalità "Login".
    Su questa pagina, Mario inserisce le sue credenziali: username: marrossi, password: MyPassWord11 e conferma l'accesso.
    L'autenticazione va a buon fine e il cliente viene reindirizzato alla home page del sito.

    \subsection*{Sxx: Modifica Password}
    \addcontentsline{toc}{subsection}{\protect\numberline{}Sxx: Modifica Password}
    Il cliente Luca Neri è autenticato e si trova sulla homepage del sito web e desidera cambiare la sua password.
    Su questa pagina, Luca seleziona il pulsante "account".
    A questo punto, si apre un menu a tendina dal quale il cliente sceglie l'opzione "Visualizza Account".
    Il sistema reindirizza Luca alla pagina del suo account, dove sono visualizzati i dettagli del suo profilo
    (Nome: Luca, Cognome: Neri, Data di nascita: 27/05/1998, Codice Fiscale: LCUNRE98E27F205X, Username: nlucas98,
     Indirizzo e-mail: lneri98@gmail.com e Password: LaMiaPasswordSafe1!).
    Sulla pagina del suo account, il cliente individua e seleziona l'opzione "Modifica Dati Account".
    Viene presentato un modulo di modifica dati account che include un campo per la modifica della password.
    Il cliente inserisce la sua nuova password: Luc1Gatto23 nel campo "Nuova Password" e la conferma nel campo "Conferma Password".
    Il sistema verifica la validità della nuova password in base ai criteri di sicurezza stabiliti.
    A questo punto conferma la modifica e visualizza un messaggio che lo informa del successo dell'operazione.

    \subsection*{Sxx: Ricerca prodotti}
    \addcontentsline{toc}{subsection}{\protect\numberline{}Sxx: Ricerca prodotti}

    Il cliente Francesco accede alla homepage iSnapGaming, dove troviamo una vetrina di videogiochi in vendita.
    Francesco ha un obiettivo specifico: trovare un gioco di simulazione. Per fare ciò si dirige verso il filtro di
    ricerca. Una volta attivato il filtro, il sistema ci mostra una lista di categorie tra cui scegliere. Francesco
    seleziona la categoria "simulazione" e conferma la scelta premendo il pulsante "Applica filtro". Il sistema risponde
    immediatamente, aggiornando la pagina e mostrando un elenco di giochi appartenenti alla categoria "simulazione".
    Il sistema ha mostrato troppi videogames e Francesco non riesce a individuare il gioco che desidera tra le opzioni
    visualizzate. A questo punto, decide di utilizzare la funzione di ricerca dell'applicazione. Questa funzione di
    ricerca è facilmente accessibile tramite la barra di navigazione, ed è qui che digita il nome del gioco
    desiderato: "GT7". Con un semplice tocco su "Cerca", il sistema elabora la richiesta e presenta i risultati.
    Tra questi, individuiamo il gioco "GT7", visualizzandone l'immagine di copertina, una breve descrizione e il prezzo.
    Curioso di saperne di più sul gioco, Francesco clicca sulla sua immagine di copertina. In risposta, il sistema
    mostra ulteriori dettagli, tra cui l'età consigliata, le dimensioni del prodotto, la data di uscita e molto altro.
    Dopo aver esaminato attentamente le informazioni, decide di procedere all'acquisto e aggiunge il gioco al suo
    carrello virtuale utilizzando la funzionalità  "Aggiungi al carrello". Tuttavia, la sessione di shopping non finisce
    qui. Soddisfatto dell'acquisto, Francesco decide di tornare alla homepage per esplorare ulteriori prodotti e vedere
    cosa altro l'applicazione ha da offrire.

    \subsection*{Sxx: Visualizzazione prodotti}
    \addcontentsline{toc}{subsection}{\protect\numberline{}Sxx: Visualizzazione prodotti}
    Francesco effettua l'accesso al sistema con autenticazione inserisce  le sue credenziali d'accesso, username :
     "francesco@libero.it" e password "aaaabbb2" ed effettua il "Login". Il sistema rileva un'eccezione, segnalando con un
    messaggio di errore "utente non riconosciuto" e reindirizza Francesco alla pagina d'autenticazione. Francesco ripete
    il processo di autenticazione con meticolosità, prestando particolare attenzione alle specificità maiuscole/minuscole
    nella password (modificandola in "aaaAbbb2"). Questa volta, l'autenticazione ha esito positivo, il sistema riconosce
    Francesco come utente valido e lo indirizza alla homepage di iSnapGaming, dove lo accoglie con un messaggio di benvenuto.
    Francesco seleziona la funzionalità "I Miei Ordini", selezionando l'appropriato pulsante e in risposta, il sistema
    fornisce tabella che espone tutti gli ordini registrati nel sistema. Ogni riga della tabella contiene il numero
    d'ordine, una breve descrizione, il costo totale e lo stato corrente dell'ordine.
    \begin{itemize}
        \item [-] \#0567 Ghostrunner 2 	\officialeuro 40.98	consegnato
        \item [-] \#1678 GT7... \officialeuro 95.49	in lavorazione
    \end{itemize}
    Francesco seleziona l'ordine \#1678 per esaminare in dettaglio l'ordine ed il  sistema visualizza una seconda tabella
    che mostra quantità, nomi e prezzi di ciascun articolo dell'ordine:
    \begin{itemize}
        \setlength\itemsep{.05cm}
        \item [-] 1	GT7 \officialeuro 35.50
        \item [-] 1Assassin's Creed Mirage	\officialeuro 59.99
    \end{itemize}
    Alla conclusione della sua sessione d'interazione con iSnapGames, Francesco seleziona la funzione
    "Logout". In risposta, il sistema lo congeda riportandolo alla homepage generica di iSnapGaming.

    \subsection*{Sxx: Aggiornamento stato ordine}
    \addcontentsline{toc}{subsection}{\protect\numberline{}Sxx: Aggiornamento stato ordine}
    Luca Rossi, il gestore degli ordini, si autentica con successo nel sistema utilizzando le proprie
    credenziali di accesso. Il sistema lo reindirizza automaticamente alla sezione dedicata alla gestione
    degli ordini, e elenca tutti gli ordini registrati.
    Al fine di semplificare il suo lavoro e focalizzarsi sugli ordini "pronti", Luca decide di applicare
    un filtro tramite l'apposita funzione. La tabella degli ordini viene quindi prontamente aggiornata per
    visualizzare solo gli ordini con lo stato "pronto". Tra questi ordini, Luca individua l'ordine con numero
    identificativo \#1234, che è stato recentemente spedito. Di conseguenza, Luca modifica lo stato di
    quest'ordine da "pronto" a "spedito". Luca prosegue il suo lavoro gestionale, aggiornando con successo
    altri cinque ordini, garantendo così che le informazioni siano sempre aggiornate e accurate.
    Successivamente, utilizza nuovamente la funzione di filtro per selezionare gli ordini con stato "spedito".
    Luca riceve la conferma della consegna dell'ordine \#4321, il quale transita dallo stato "spedito" a
    "consegnato". La giornata di lavoro di Luca continua con ulteriori aggiornamenti degli ordini,
    garantendo che ciascun ordine possa passare attraverso i vari stati possibili, tra cui "in preparazione",
    "pronto", "spedito" e "consegnato."

    \subsection*{Sxx: Gestione account}
    \addcontentsline{toc}{subsection}{\protect\numberline{}Sxx: Gestione account}
    Il cliente Francesco a seguito di un cambio di residenza decide di aggiornare le informazioni su
    iSnapGaming, effettua l'accesso al sistema tramite autenticazione e  seleziona la funzionalità
    "modifica Account". Il sistema mostra una form con tutti i dati (nome,cognome, città di residenza,
    provincia, via, numero civico), Francesco modifica città : Roma, provincia : RO, via: viale Italia,
    numero civico : 1 e utilizza la funzionalità "aggiorna" per rendere i dati persistenti.

    %COPIARE ED INCOLLARE IL SEGUENTE BLOCCO DI ISTRUZIONI - INSERIRE NOME SCENARIO NEI DUE APPOSITI CAMPI.
    \subsection*{Sxx: INSERISCI NOME SCENARIO}
    \addcontentsline{toc}{subsection}{\protect\numberline{}Sxx: INSERISCI NOME SCENARIO}
        Descrivi qui il tuo scenario.




\end{document}